\documentclass[11pt]{article}
\usepackage[a4paper,margin=25mm]{geometry}
\usepackage{amsmath,amssymb,amsthm,mathtools,booktabs,longtable}
\usepackage[hidelinks]{hyperref}
\newtheorem{theorem}{Theorem}
\newtheorem{lemma}[theorem]{Lemma}
\newtheorem{proposition}[theorem]{Proposition}
\newcommand{\C}{\mathbb C}
\newcommand{\R}{\mathbb R}
\newcommand{\im}{\operatorname{im}}
\newcommand{\diag}{\operatorname{diag}}
\newcommand{\rank}{\operatorname{rank}}
\newcommand{\bdeg}{\operatorname{bideg}}
\newcommand{\ord}{\operatorname{ord}}
\title{Vanishing of the original proper-source kernel\ through conductor coprimality}
\author{Split-Zero programme: native continuation}
\date{19 September 2026}
\begin{document}
\maketitle
\begin{abstract}
The supplied native continuation proves that the actual proper-source
kernel is annihilated by the fourth power of the product of the four
Fourier coordinates. We evaluate the remaining coprimality question in
the original period family. For every retained arithmetic unit phase,
the limiting conductor and coordinate divisor are coprime on the domain
$\gamma\ge7\delta>0$. An explicit family of sixteen $10$ by $6$
coefficient matrices extends this to a finite large-period domain and
to the complement of a finite exceptional set in the admitted period
domain. Consequently the original proper-source kernel is zero there,
at every original degree $k\ge9$. The source Schur term disappears by
an equality of the original spaces. The full target minimum remains in
the metric receiving this result. All coordinates and unit factors are
restored by explicit maps.
\end{abstract}

\section{Original objects and the question settled here}
Let $0<\delta<1/2$, $\gamma>2$, and retain the ordered simple quartet
\[
 w_1=\delta+i\gamma,\quad w_2=\delta-i\gamma,
 \quad w_3=-\delta+i\gamma,\quad w_4=-\delta-i\gamma.
 \tag{NC1}
\]
Put $\beta=2(\gamma^2-\delta^2)$ and
$\eta=(\delta^2+\gamma^2)^2$. The centered polynomial is
$h_c(w)=w^4+\beta w^2+\eta$. The evaluation matrix $V$ has rows
$(1,w_j,w_j^2,w_j^3)$. Retain the complete unit matrix
$U=\diag(a,\bar a,\bar a,a)$, where the actual $a\ne0$ is fixed.
The original period family from PCL1--7 \cite{PCL} is
\[
 H(u)=D_u\mathcal R(u^{-1})V^{-1}U,
 \qquad D_u=e^{\mathfrak a/u}G_u,
 \quad\mathfrak a=1/160+\beta/24+\eta/2. \tag{NC2}
\]
Here $G_u$ is the original invertible diagonal matrix of Gamma ray
constants and branch powers. Its values are retained. The entire matrix
$\mathcal R$ has determinant one and
\[
 J_\infty=\mathcal R(0)=
 \begin{pmatrix}1&0&-\beta/3&0\\0&1&0&-2\beta/3\\
 0&0&1&0\\0&0&0&1\end{pmatrix}. \tag{NC3}
\]
For completeness, the full coefficients used below are
\[
 R_{br,n}=\sum_{\substack{p,h\ge0\\2p+4h=5n+r+1-b}}
 \frac{\beta^p\eta^h}{3^p p!h!}(-1)^{p+h-n}
 5^{p+h-n}(b/5)_{p+h-n},
 \quad1\le b\le4,\quad0\le r\le3. \tag{NC4}
\]
The convergence and determinant identity are the supplied PCL proof;
they are not inferred from a finite truncation. This paper proves the
additional coprimality and kernel conclusions from that exact family.

Write $D=\diag(\zeta,\zeta^2,\zeta^3,\zeta^4)$,
$\zeta=e^{2\pi i/5}$, and $g f(X)=f(D^{-1}X)$. In the literal
Fourier coordinates $X=(X_1,\ldots,X_4)$ set
\[
 Q_u(X)=Q_0(H(u)^{-1}X),\quad Q_0(t)=t_1t_4-t_2t_3,
 \quad Q_{u,j}=g^jQ_u,\quad A_u=\prod_{j=1}^4Q_{u,j}.
 \tag{NC5}
\]
Work on the original five-orbit domain: these five quadrics are
nonassociate prime polynomials. Retain the original degrees
$k\equiv1\pmod4$, $k\ge9$. Let $R_u=\C[X]/(Q_u)$,
$B_u=\im(\C[X]^{\langle g\rangle}\to R_u)$ and
$C_{u,k}=(R_u/(A_u))_k$. Let $J_{u,k}$ denote the image of invariant
degree-$k$ polynomials in $C_{u,k}$; its original first-$v$-jet-vanishing
subspace is denoted by $J_{v,k}$. The proper-source kernel under study is
exactly
\[
 E_k=\{f\in J_{v,k}: X_i^4 f\in J_{u,k+4}
                 \text{ for }i=1,2,3,4\}. \tag{NC6}
\]
All memberships here are in the specified coefficient quotients. These
are the cofactor kernel and four multipliers of the supplied note E1--4
\cite{Native}. There is no replacement of this kernel by an arithmetic
eigenspace or by the one-step observation kernel.

\section{Return from the conductor to the actual source kernel}
We give the algebraic argument, including the vertex, so that the
kernel conclusion does not depend on a rank analogy.
\begin{lemma}
The conductor $\{c\in R_u:cR_u\subset B_u\}$ equals $A_uR_u$.
With $p=X_1X_2X_3X_4$, the class $p^4 f$ vanishes in $C_{u,k+16}$
for every $f\in E_k$.
\end{lemma}
\begin{proof}
For any homogeneous lift $h$, the invariant sum
$\sum_{j=0}^4g^j(A_uh)$ restricts on $Q_u=0$ to $A_uh$:
each other summand contains $Q_u$. Thus $A_uR_u\subset B_u$,
without dividing the sum by five.

For the reverse inclusion localize at $X_i^5$. The cyclic action on
$X_i\ne0$ is free. The localized polynomial ring is free over its
invariants with basis $1,X_i,\ldots,X_i^4$: each character can be
cleared by an invertible power of $X_i$. The relation is
$T^5-X_i^5$ and its derivative is a unit. This faithfully flat base
change turns the invariant-image inclusion into
\[
 \C[X]/\Bigl(\prod_{j=0}^4Q_{u,j}\Bigr)
       \longrightarrow\bigoplus_{j=0}^4\C[X]/(Q_{u,j})
 \tag{NC7}
\]
on that chart. Indeed an invariant polynomial vanishing modulo $Q_u$
vanishes modulo every $Q_{u,j}$, hence is divisible by their product;
its quotient is again invariant. The free cyclic orbit gives (NC7).
A tuple supported on its $j$th component lies in the image exactly when
that component is divisible by all four other primes. Their ideals
intersect in their product in the polynomial UFD. Every ideal in the
direct sum decomposes componentwise, proving the conductor description
on each $X_i\ne0$ chart. Faithful flatness descends that description.

The smooth projective quadric has its original Segre coordinates, with
$(R_u)_d=T_{d,d}$, the biforms of bidegree $(d,d)$.
If a homogeneous $c$ of degree $d$ is in the conductor, $c/A_u$ is
a regular section of $\mathcal O(d-8,d-8)$ on all four $X_i$ charts.
They cover the projective quadric. On the standard Segre charts the
regularity conditions say precisely that this section is a biform of
that bidegree, or is zero if $d<8$. Thus $c\in A_uR_u$ globally.
There is no extra vertex-supported contribution.

Lift $f\in E_k$ to $B_u$. By (NC6) and $A_uR_u\subset B_u$, each
$X_i^4f$ is in $B_u$. For a monomial $m$ of nonzero character $w$,
choose $i$ with $4i=w\pmod5$. The polynomial $p^4m/X_i^4$ has
nonnegative exponents and character zero. Hence $p^4mf\in B_u$.
For a weight-zero monomial use $f\in B_u$ directly. Monomials span
$R_u$, so $p^4f$ is in the conductor and therefore in $A_uR_u$.
\end{proof}

In the original Segre substitution let $\ell_i=X_i|_{Q_u=0}$ and
$\mathcal A=A_u|_{Q_u=0}$. They have bidegrees $(1,1)$ and $(8,8)$.
The previous lemma and polynomial divisibility give the exact ambient
multiplication kernel, with all multiplicities retained:
\[
 \ker M_{p^4}=
 (\mathcal A/G)T_{k-8+g_1,k-8+g_2}/
                    \mathcal A T_{k-8,k-8},
 \quad G=\gcd(\mathcal A,(\ell_1\ell_2\ell_3\ell_4)^4),
 \quad\bdeg G=(g_1,g_2). \tag{NC8}
\]
Indeed $\mathcal A\mid p^4f$ is equivalent to
$\mathcal A/G\mid f$ in the biform UFD. Consequently $G=1$ implies
$E_k=0$. We now prove $G=1$ on an explicit original period domain.

\section{An exact auxiliary coordinate map, with its scalar retained}
Set $z=u^{-1}$ and $Y=D_u^{-1}X$. This map is invertible, commutes
with $D$, and carries every coordinate hyperplane to itself. Its exact
polynomial identities are
\[
 \begin{aligned}
 Q_u(D_uY)&=\widetilde Q_z(Y)
 :=Q_0(U^{-1}V\mathcal R(z)^{-1}Y),\\
 A_u(D_uY)&=\prod_{j=1}^4\widetilde Q_z(D^{-j}Y),
 \qquad p(D_uY)=\det(D_u)p(Y).
 \end{aligned}\tag{NC9}
\]
Thus coprimality and ideal membership are transported in both directions.
This is a coordinate isomorphism for a divisibility calculation; all
$D_u$ factors stay in the original Hermitian metrics.

Write
\[
 t=\frac{\beta}{6}=\frac{\gamma^2-\delta^2}{3}>0,
 \quad d=2\delta\gamma>0,
 \quad\alpha=(\delta+i\gamma)^2=-3t+id,
 \quad c=a^{-2}\ne0. \tag{NC10}
\]
At $z=0$, (NC3) gives the coefficient polynomial
\[
 P_Y(w)=(Y_1+2tY_3)+(Y_2+4tY_4)w+Y_3w^2+Y_4w^3.
 \tag{NC11}
\]
The complete original quadric is consequently
\[
 \widetilde Q_0(Y)=
 c\{[Y_1+(\alpha+2t)Y_3]^2
       -\alpha[Y_2+(\alpha+4t)Y_4]^2\}
 -\bar c\{[Y_1+(\bar\alpha+2t)Y_3]^2
       -\bar\alpha[Y_2+(\bar\alpha+4t)Y_4]^2\}.
 \tag{NC12}
\]
The subtraction is coefficientwise; it does not conjugate the variables.
Introduce the real-coefficient polynomial $q_0$ by the exact identity
$\widetilde Q_0=2i q_0$. Thus its four-factor conductor has the scalar
$(2i)^4$ as well. In the order $(Y_1,Y_3;Y_2,Y_4)$, write
\[
 q_0(Y)= (Y_1,Y_3)A(Y_1,Y_3)^T
             +(Y_2,Y_4)B(Y_2,Y_4)^T. \tag{NC13}
\]
The real symmetric blocks are specified without any unit-phase choice:
\[
 A=\begin{pmatrix}\Im c&\Im(c(-t+id))\\
 \Im(c(-t+id))&\Im(c(-t+id)^2)\end{pmatrix},\qquad
 B=-\begin{pmatrix}\Im(c\alpha)&\Im(c\alpha(t+id))\\
 \Im(c\alpha(t+id))&\Im(c\alpha(t+id)^2)\end{pmatrix}.
 \tag{NC14}
\]
For complex $C$ and $L$, direct expansion yields
\[
 \Im C\,\Im(CL^2)-\Im(CL)^2=-|C|^2(\Im L)^2.
\]
It follows that
\[
 \det A=-|c|^2d^2<0,\qquad
 \det B=-|c|^2|\alpha|^2d^2<0. \tag{NC15}
\]
In particular both real binary quadratic forms are nonsingular and
indefinite. No special unit value has been imposed.

\section{All coordinate-plane intersections at the original limit}
\begin{lemma}
If $0<d\le t$, no diagonal entry of $A$ and diagonal entry of $B$
can both be zero. Neither pair of diagonal entries of a single block
can both be zero.
\end{lemma}
\begin{proof}
The four diagonal coefficients are $\Im(cf_j)$ for
\[
 f_1=1,\quad f_2=(-t+id)^2,\quad
 f_3=-\alpha,\quad f_4=-\alpha(t+id)^2.
\]
If $\Im(cf_j)=\Im(cf_l)=0$, then $f_j\bar f_l$ is real. The exact
cross-block values are
\[
\begin{array}{c|cc}
 \Im(f_j\bar f_l)&l=3&l=4\\\hline
 j=1&d&-d(d^2+5t^2)\\
 j=2&-d(d^2+5t^2)&d(d^4+6d^2t^2-11t^4)
\end{array} \tag{NC16}
\]
The bottom right entry is negative for $0<d\le t$, since its
parenthesis is at most $-4t^4$. The other entries are nonzero.
Within a block, the ratios $f_2/f_1=(-t+id)^2$ and
$f_4/f_3=(t+id)^2$ have nonzero imaginary parts. These observations
prove both assertions, for every nonzero $c$.
\end{proof}

\begin{theorem}
If $0<d\le t$, then for every $i=1,2,3,4$ and $j=1,2,3,4$ the
two ternary quadrics
\[
 q_0|_{Y_i=0},\qquad q_0(D^{-j}Y)|_{Y_i=0}
 \tag{NC17}
\]
are coprime. Therefore the limiting conductor is coprime to the
product of the four coordinate sections on the original quadric.
\end{theorem}
\begin{proof}
Deleting one variable from a block in (NC13) leaves a single square
in that block and the other nonsingular binary form. If the square
coefficient is nonzero, the ternary form has rank three, hence is an
irreducible conic over $\C$. A common factor with its diagonal
transform would then force the two conics to be proportional.
For $j\ne0\pmod5$, proportionality of a polynomial and its transform
requires every nonzero monomial to have the same weight modulo five.
The weights of the six possible monomials are
\[
\begin{array}{c|rrrrrr}
 \text{monomial}&Y_1^2&Y_1Y_3&Y_3^2&Y_2^2&Y_2Y_4&Y_4^2\\
 \text{weight}&2&4&1&4&1&3
\end{array}. \tag{NC18}
\]
On $Y_2=0$ or $Y_3=0$ all remaining monomial weights are distinct;
rank three excludes a one-monomial polynomial. On $Y_1=0$, the only
weight supporting two monomials is weight one, requiring both diagonal
entries of $B$ to vanish. On $Y_4=0$, it is weight four, requiring both
diagonal entries of $A$ to vanish. The preceding lemma excludes both.
Thus the irreducible case is settled.

If the square coefficient is zero, the ternary conic is the other
binary form. That form has two distinct real projective roots by
(NC15). Both are finite and nonzero: a zero or infinite root would
make a diagonal entry of that block zero, which (NC16) forbids
simultaneously with the deleted block's diagonal coefficient.
The ratio of its two coordinates is multiplied under $D^{-j}$ by
$\zeta^{2j}$ or its inverse. For $1\le j\le4$ this number is
nonreal. Multiplying a finite nonzero real root by it cannot yield
either real root. Hence the two reducible conics share no line.
These cases exhaust every coordinate-plane section and every retained
unit phase.

A common divisor of a coordinate section and the conductor on the
Segre surface would give a curve on $Q_0=Q_j=Y_i=0$ for some $i,j$.
In the plane $Y_i=0$ this is precisely a common polynomial factor in
(NC17), just excluded. This proves the last assertion.
\end{proof}

The condition is established for every actual hypothetical off-line
quartet: $\gamma\ge7\delta$ gives
$d/t=6\delta\gamma/(\gamma^2-\delta^2)\le7/8<1$.
Platt and Trudgian's rigorously verified finite-height theorem
\cite{PT} places any such zero above $3\cdot10^{12}$, whereas
$\delta<1/2$. Only the much weaker consequence $\gamma\ge7\delta$
is used here. The cited theorem is a human analytic input; the
coprimality argument above is proved independently for its full stated
parameter domain. For a formal non-zeta quartet outside that domain,
the finite matrices below still decide coprimality, but the preceding
universal nonvanishing argument is not asserted there.

\section{Sixteen finite matrices give an original finite-period bound}
For each $i,j$ let $q_i(z)$ and $q_{ij}(z)$ be the restrictions to
$Y_i=0$ of $\widetilde Q_z(Y)$ and $\widetilde Q_z(D^{-j}Y)$.
They are ternary quadrics with entire coefficient functions. In the
unchanged ordered monomial bases define the $10$ by $6$ matrix
\[
 M_{ij}(z):(L_1,L_2)\longmapsto L_1q_i(z)+L_2q_{ij}(z),
 \qquad (L_1,L_2)\in T_1(\C^3)^2,\quad
 M_{ij}(z)(L_1,L_2)\in T_3(\C^3). \tag{NC19}
\]
Here $T_e(\C^3)$ means the homogeneous ternary forms of degree $e$;
it is not the biform notation used earlier. Neither restricted
quadric is zero, since an invertible transform of $Q_0$ cannot have
a coordinate hyperplane as a factor.

\begin{lemma}
The matrix $M_{ij}(z)$ has rank six if and only if its two quadrics
are coprime.
\end{lemma}
\begin{proof}
If they are coprime, $L_1q_i=-L_2q_{ij}$ implies $q_i\mid L_2$ in
the polynomial UFD. Degree forces $L_2=0$, then $L_1=0$.
If their gcd has degree one, write them as $g a,g b$ with $a,b$
linear. The pair $(b,-a)$ is a nonzero kernel vector. If the gcd has
degree two, the quadrics are proportional and multiplication of that
constant relation by any linear form gives a kernel vector.
\end{proof}

Expand $M_{ij}(z)=\sum_{n\ge0}M_{ij,n}z^n$, using exactly (NC4),
the polynomial adjugate $\mathcal R^{-1}=\operatorname{adj}\mathcal R$,
and (NC9). In Euclidean coefficient norms set
\[
 \mu_{ij}=\sigma_{\min}(M_{ij,0})>0,\qquad
 C_{ij}=\sum_{n\ge1}\|M_{ij,n}\|_2<\infty,
 \qquad
 R_{\rm cop}=\max\left(1,\max_{i,j}\frac{2C_{ij}}{\mu_{ij}}\right).
 \tag{NC20}
\]
The just-proved rank theorem establishes each strict positivity; it is
not an assumed nonzero minor. Entire convergence gives the summability
in (NC20), for example by a Cauchy bound on the circle of radius two.
All coefficients in this finite list of specified series are given by
the original packet, its actual unit and (NC4). For $|u|\ge R_{\rm cop}$,
\[
 \|M_{ij}(u^{-1})-M_{ij,0}\|_2\le C_{ij}/|u|
 \le\mu_{ij}/2,
 \qquad \sigma_{\min}(M_{ij}(u^{-1}))\ge\mu_{ij}/2>0.
 \tag{NC21}
\]
This proves coprimality at finite original periods. If $C_{ij}=0$ its
matrix is constant and the same conclusion follows immediately.

One can also retain the whole original domain $|u|\ge R_*>0$.
For each matrix choose the first ordered nonzero $6$ by $6$ minor at
$z=0$, and let $f_{ij}(z)$ denote that same minor at $z$. Their
product $F(z)$ is entire and $F(0)\ne0$. Its zeros in
$|z|\le R_*^{-1}$ are finite. Outside that explicitly specified finite
set every matrix has rank six. At a zero of $F$, use all minors of
(NC19): a zero of the selected minor alone does not assert failure.

\begin{theorem}
For an actual hypothetical simple off-line quartet, on the original
five-orbit domain and every fixed period $|u|\ge R_{\rm cop}$,
\[
 \gcd\bigl(\mathcal A,(\ell_1\ell_2\ell_3\ell_4)^4\bigr)=1,
 \qquad E_k=0\quad(k\ge9). \tag{NC22}
\]
The same conclusion holds at every admitted period outside the finite
exceptional set specified above. Period, branch and unit are fixed
while $k$ grows.
\end{theorem}
\begin{proof}
Apply (NC21), the rank criterion, the coordinate-plane argument, the
exact pullback (NC9), and (NC8), in that order. A divisor occurring
with higher multiplicity would in particular occur once, already
excluded. This retains every factor multiplicity of the original
conductor. The admissible five-orbit radius, if larger, is imposed by
taking the maximum with $R_{\rm cop}$.
\end{proof}

\section{Return to the original rank and attained metric}
The supplied low-jet calculation J1--6 \cite{Native} defines
$v=\ord_0E_A$ at the same actual period. On its two-active-coordinate
locus it proves $t_0(k,u)=0$ for $k\ge5v-1$. It also proves that the
one-coordinate exceptions in the admitted inverse-period disk lie in
a finite set. Intersecting this locus with (NC22) therefore gives
\[
 j_{0,k}=8k-32-v,\qquad r_I=j_{0,k}-\dim E_k=8k-32-v
 \quad\bigl(k\ge\max(9,5v-1)\bigr). \tag{NC23}
\]
At all other periods the exact identity is still
$r_I=j_{0,k}-\dim E_k$ with the actual finite jet rank and conductor
gcd; the displayed value of $t_0$ is not substituted there.

Here is the actual metric consequence. Use the original notation of
the proper-source map: $Z_R$ is its ordered value-space inclusion,
$J_0$ its coefficient lift, and $D_{k,N-v}^{(s_0)}$ the original
degree-$k$, source-order-$s_0$ attained form at its original cutoff.
Before (NC22), this source metric is the attained quotient by $E_k$:
for an $E_k$ basis matrix $Z_E$ it is
\[
 Q_N^p=Z_R^*J_0^*D J_0 Z_R
 -Z_R^*J_0^*D J_0 Z_E
 (Z_E^*J_0^*D J_0 Z_E)^{-1}Z_E^*J_0^*D J_0 Z_R,
 \quad D=D_{k,N-v}^{(s_0)}. \tag{NC24}
\]
This is obtained by completing the square in the actual affine fiber;
the second term uses the full cross block. On (NC22) the fiber
direction space is zero. Consequently the exact receiving equality is
\[
 Q_N^p=Z_R^*J_0^*D_{k,N-v}^{(s_0)}J_0Z_R. \tag{NC25}
\]
The target metric retains its entire target direction space
$\mathcal J$. If $W$ is its actual value lift, $T$ a basis of
$\mathcal J$, and $G$ the original target form, completing the square
still gives
\[
 H_N^p=W^*GW-W^*GT(T^*GT)^{-1}T^*GW. \tag{NC26}
\]
No direction in $\mathcal J$ has been removed by (NC22). The original
relative determinant return is thus exactly
\[
 \mathcal M^p=\mathcal R_{\rm end}\log\det
 \left[(Z_R^*J_0^*D J_0 Z_R)^{-1/2}
        H_N^p(Z_R^*J_0^*D J_0 Z_R)^{-1/2}\right]. \tag{NC27}
\]
$\mathcal R_{\rm end}$ applies the original four endpoint signs.
Equation (NC27) removes an actual source minimization. Evaluation of
its remaining target minimum and native relative spectrum requires
the original matrices in (NC26); the ambient displacement coefficient
in D24 is not inserted as their determinant.

\section{Backward receiving statements and the next calculation}
The following replacements are precise consequences for the supplied
four notes. In E4, the coefficient $(g_z+g_w)(k-7)+g_zg_w$ is zero
on (NC22), and $E_k$ itself is zero. In its final metric paragraph,
the previously unevaluated coprime stratum is now proved to contain
the finite-period domain (NC20)--(NC21), for every actual retained unit
phase. J6 keeps the original observation kernel and is unchanged;
(NC23) concerns the independent proper-source rank. O9--12 continue
to recover the original full packet by observation iterates; none of
their inverses has been substituted into the target minimum (NC26).

The scalar results D19, D20, D23 and D24 retain their original
four-endpoint domains and analytic providers. The two supplied
auxiliary scripts check their displayed coefficient identities and
character counts (11 and 45 checks); those checks do not constitute
another proof of the imported analytic estimates. The four complete
incoming proof notes accompany this continuation. The separately
supplied low-jet note adds the smaller one-coordinate matrix J3a and
explicit classes J3b and is retained as the newer version.

The next numerical or asymptotic metric calculation now starts with
the literal positive pair in (NC25)--(NC26), at the proper-source
displaced cutoffs, with no source $E_k$ block. This is the map on
which the remaining relative eigenvalues and determinant must be
computed. The theorem does not assign individual projected signs or
the seven absolute allocations from this rank result.

\begin{thebibliography}{9}
\bibitem{Native} \emph{Programme Native Continuation}, supplied
19 September 2026: ``Displaced lower covariance and matched native
return through order $q$'' (D1--25); ``Original observation: cofinal
low-degree kernel ranks'' (J1--6, including J3a--b in the separately
supplied revision); ``The original proper-source kernel is annihilated
on the coordinate divisor'' (E1--4); ``Full original packet recovery
from iterated observation'' (O1--12). Complete source files are
included with this continuation, with their original bytes preserved.
\bibitem{PCL} \emph{The literal original conductor limit and an
invertible pivot restriction}, PCL1--7 and PCL2's complete coefficient
series, existing programme proof. Included as
\texttt{PCL\_COMPLETE.tex}. The scalar, diagonal, branch,
Vandermonde and arithmetic unit factors are those of PCL7.
\bibitem{PT} Dave Platt and Tim Trudgian,
\emph{The Riemann hypothesis is true up to $3\cdot10^{12}$},
Bulletin of the London Mathematical Society \textbf{53} (2021),
792--797. \href{https://doi.org/10.1112/blms.12460}{doi:10.1112/blms.12460};
\href{https://arxiv.org/abs/2004.09765}{arXiv:2004.09765}.
Used only for the proved finite-height exclusion of off-line zeros in
the parameter-domain check following (NC18).
\end{thebibliography}
\end{document}
